
\section{Introduction}
\label{sec:introduction}
\noindent
\enquote{A major limiting factor to the widespread use of information visualization is the difficulty of labeling information abundant displays}~\cite{Fekete1999}.
This assessment of Fekete and Plaisant is as tenable today as it was in 1999.
Most of today's visualizations are still troubled by missing labels, leading to one-by-one assessment of graphical elements, overlapping labels obscuring depicted data, and ambiguous association of labels to graphical elements.
Even recent advances in big-data visualization have not managed to resolve these well-known labeling issues, despite their tremendous potential for increasing the visualization quality. 

We suppose there are at least two reasons for labeling to remain mostly unattended in real-time computer graphics.
First, high-quality text rendering and placement is an intricate task, especially in 3D. 
It starts with text rendering itself requiring lots of implementation and use-case specific tweaking and probably ends with having no text rendering support in any of the latest rendering APIs, be it OpenGL~\cite{opengl}, WebGL~\cite{webgl}, or Vulkan~\cite{vulkan}.
In research, this means lots of available resources ought to be allocated for prototype development rather than novelty and evaluation.
Second, other industries, like the gaming industry usually get along with basic (bitmap based) text rendering.
With a few exceptions (e.g., \emph{Tom Clancy's The Division}, developed by Massive Entertainment and published by Ubisoft in 2016, shows impressive dynamic labeling in its UI and world map), high-end labeling is attributed to artistic rather than programming efforts and is the result of tweaking and pre-computing assets in compliance with a title's art direction.
Furthermore, the number of labels is typically much smaller and sophisticated placement requirements are not present. 
Accordingly no industry seems to drive the progression or standardization of labeling technology in general.

In order to initiate a research driven specification of an open label library, we (1) put our idea of a unified labeling API forward for discussion, and (2) propose subsequent steps required for a widespread use in information visualization.
To this extent, the following contributions are made: %
\begin{itemize}
	\item a brief introduction to and overview of existing text rendering approaches and systems is provided,
	\item snippets of the labeling API \emph{OpenLL}, derived from use cases of our research group, are presented, and
	\item our ongoing open-source implementation efforts for mobile and desktop applications are described.
\end{itemize}

\smallskip
\noindent
OpenLL is designed to complement existing rendering APIs, systems, and transmission formats (such as glTF~\cite{gltf}).
It facilitates a descriptive labeling approach using label and proxy objects that are processed by means of renderers that are aligned to implementation specifics of hardware-accelerated text rendering techniques.
Thus, this paper provides the general concepts of OpenLL instead of drafting an extension to any rendering API or transmission format.
